% ===================================================================
% Section 5: MRAC 
% ===================================================================

\section{MRAC regulácia – gradientná metóda}
\label{sec:mrac_gradient}

V tejto kapitole je riešená adaptívna regulácia typu
MRAC (Model Reference Adaptive Control),
ktorá umožňuje automatické prispôsobovanie parametrov regulátora
pri nepresne známej dynamike riadeného systému.
Na rozdiel od klasických regulátorov s pevnými parametrami
MRAC využíva referenčný model,
ktorý definuje požadované správanie uzavretého regulačného obvodu,
a adaptačný mechanizmus,
ktorý zabezpečuje sledovanie tohto správania v čase.

V ďalšej časti je popísaný samotný algoritmus gradientnej MRAC regulácie,
vrátane voľby referenčného modelu,
štruktúry zákona riadenia
a adaptačného zákona založeného na gradientnej metóde.

% -------------------------------------------------------------------
% Subsection - Popis algoritmu
% -------------------------------------------------------------------
\subsection{Popis algoritmu}
\label{sec:mrac_gradient_algorithm}

Gradientná MRAC regulácia je založená na myšlienke,
že dynamika riadeného systému má byť v čase prispôsobovaná tak,
aby výstup systému sledoval výstup zvoleného referenčného modelu.
Návrh algoritmu pozostáva z troch základných častí:
voľba referenčného modelu,
návrh zákona riadenia
a odvodenie adaptačného zákona.

\paragraph{Model riadeného systému}

Uvažujme lineárny spojitý systém druhého rádu
\begin{equation}
	G_s(s) = \frac{b_0}{s^2 + a_1 s + a_0},
	\label{eq:mrac_plant}
\end{equation}
kde $a_0$, $a_1$, $b_0$ sú neznáme konštanty.
Predpokladá sa znalosť znamienka parametra $b_0$.

\paragraph{Referenčný model a statické zosilnenie}

Pri voľbe referenčného modelu je potrebné zabezpečiť,
aby jeho statické zosilnenie bolo rovné jednej,
t.\,j.
\[
K = 1,
\]
aby nedochádzalo k nechcenému zosilňovaniu referenčného signálu.

Pre systém druhého rádu zvolíme referenčný model
\begin{equation}
	\frac{y_m(s)}{r(s)} =
	\frac{b_{0m}}{s^2 + a_{1m}s + a_{0m}},
	\label{eq:mrac_ref_model}
\end{equation}
pričom sa volí
\begin{equation}
	b_{0m} = a_{0m},
\end{equation}
čím sa zabezpečí jednotkové statické zosilnenie modelu.

\paragraph{Zákon riadenia}

Všeobecný tvar zákona riadenia pri splnenej podmienke zhody je
\begin{equation}
	u(t) = \Theta^T(t)\,\omega(t),
	\label{eq:mrac_general_control}
\end{equation}
kde regresný vektor obsahuje merateľné signály.

Pre systém druhého rádu zvolíme
\begin{equation}
	\omega(t) =
	\begin{bmatrix}
		y(t) \\
		\dot{y}(t) \\
		r(t)
	\end{bmatrix}.
\end{equation}

Zákon riadenia potom nadobúda tvar
\begin{equation}
	u(t) = \Theta_1 y(t) + \Theta_2 \dot{y}(t) + \Theta_3 r(t).
	\label{eq:mrac_control_law}
\end{equation}

\paragraph{Podmienka zhody a ideálne parametre}

Po dosadení \eqref{eq:mrac_control_law} do systému
\eqref{eq:mrac_plant} dostávame podmienku zhody

\begin{equation}
	y(s)\left(
	s^2 + (a_1 - b_0\Theta_2)s + a_0 - b_0\Theta_1
	\right)
	=
	b_0\Theta_3 r(s).
\end{equation}

Prenos uzavretého regulačného obvodu je
\begin{equation}
	\frac{y(s)}{r(s)} =
	\frac{b_0\Theta_3}
	{s^2 + (a_1 - b_0\Theta_2)s + (a_0 - b_0\Theta_1)}.
\end{equation}

Porovnaním s referenčným modelom získame ideálny vektor parametrov

\begin{align}
	\Theta_1^\ast &= \frac{a_0 - a_{0m}}{b_0}, \\
	\Theta_2^\ast &= \frac{a_1 - a_{1m}}{b_0}, \\
	\Theta_3^\ast &= \frac{b_{0m}}{b_0}.
\end{align}

\paragraph{Adaptačná odchýlka}

Adaptačná odchýlka je definovaná ako
\begin{equation}
	e(t) = y(t) - y_m(t),
\end{equation}
pričom cieľom regulácie je dosiahnuť
\begin{equation}
	\lim_{t \to \infty} e(t) = 0.
\end{equation}

\paragraph{Účelová funkcia a MIT pravidlo}

Minimalizovaná je účelová funkcia
\begin{equation}
	J(\Theta) = \frac{1}{2} e^2(t).
\end{equation}

MIT pravidlo je formulované ako
\begin{equation}
	\frac{d\Theta}{dt}
	=
	-\alpha \frac{\partial J}{\partial \Theta}
	=
	-\alpha e \frac{\partial e}{\partial \Theta}.
\end{equation}

Keďže $y_m$ nezávisí od $\Theta$, platí
\begin{equation}
	\frac{d\Theta}{dt}
	=
	-\alpha e \frac{\partial y}{\partial \Theta}.
\end{equation}

\paragraph{Zákon adaptácie pre systém druhého rádu}

Pre systém druhého rádu sú adaptačné zákony dané vzťahmi

\begin{align}
	\dot{\Theta}_1
	&=
	-\alpha_1 e \,
	\mathrm{sign}(b_0)
	\frac{1}{s^2 + a_{1m}s + a_{0m}}\,y,
	\\
	\dot{\Theta}_2
	&=
	-\alpha_2 e \,
	\mathrm{sign}(b_0)
	\frac{s}{s^2 + a_{1m}s + a_{0m}}\,y,
	\\
	\dot{\Theta}_3
	&=
	-\alpha_3 e \,
	\mathrm{sign}(b_0)
	\frac{1}{s^2 + a_{1m}s + a_{0m}}\,r.
\end{align}

Člen $s$ v adaptačnom zákone pre parameter $\Theta_2$
zohľadňuje dynamický vplyv derivácie výstupu
v zákone riadenia.

% -------------------------------------------------------------------
% Subsubsection - Diskrétna realizácia steady-state
% -------------------------------------------------------------------
\subsubsection{Diskrétna realizácia (real-time implementácia)}
\label{sec:mrac_gradient_discrete}

Pre praktickú implementáciu gradientnej MRAC regulácie v reálnom čase
je potrebné previesť spojitý algoritmus do diskrétneho tvaru tak,
aby v každom vzorkovacom kroku boli použité iba aktuálne merané veličiny
a stavové premenné uložené z predchádzajúceho kroku.

Predpokladá sa, že zo systému je dostupný iba výstup $y(k)$,
pričom jeho derivácia je aproximovaná numericky pomocou spätnej diferencie

\begin{equation}
\dot{y}(k)
=
\frac{y(k)-y(k-1)}{T_s},
\label{eq:mrac_dy_discrete}
\end{equation}

kde $T_s$ je vzorkovací čas.



% -------------------------------------------------------------------
% Subsection - Implementácia MRAC v prostredí MATLAB
% -------------------------------------------------------------------
\subsection{Implementácia MRAC v prostredí MATLAB}
\label{sec:mrac_matlab_implementation}

Cieľom tejto časti bolo implementovať navrhnutý gradientný MRAC algoritmus
v prostredí MATLAB/Simulink bez využitia štandardných blokov riadenia,
preddefinovaných adaptívnych štruktúr alebo iných pokročilých nástrojov
Simulinku. Celý algoritmus je realizovaný výlučne v rámci jedného bloku
MATLAB Function, pričom všetky výpočty – riadiaci zákon,
referenčný model, filtračné členy aj adaptačný mechanizmus –
sú explicitne implementované v zdrojovom kóde.

Použitý blok MATLAB Function nemá vnútornú pamäť medzi
jednotlivými vzorkovacími krokmi. Z tohto dôvodu je potrebné zabezpečiť,
aby všetky stavové premenné, ktoré musia byť prenesené do nasledujúceho
diskrétneho kroku, boli vedené ako vstupno-výstupné premenné bloku.
Konkrétne ide o:

\begin{itemize}
    \item stav referenčného modelu $x_m$,
    \item stav filtra výstupu $x_{yf}$,
    \item stav filtra referencie $x_{rf}$,
    \item adaptačný vektor parametrov $\Theta$,
    \item predchádzajúcu hodnotu výstupu systému $y_{prev}$.
\end{itemize}

Tieto veličiny sú na výstupe bloku aktualizované a následne spätne
privedené na jeho vstup, čím sa zabezpečí diskrétna dynamická realizácia
celého algoritmu.

Schéma prepojenia blokov v prostredí Simulink je znázornená nižšie.

\begin{center}
    \vbox{
        \makebox[\textwidth][c]{%
            \includegraphics[width=1.3\textwidth, trim=100 50 100 50, clip]{TS_mrac_simulink_scheme.pdf}
        }
        \figcaption{Zapojenie MRAC regulátora v bloku MATLAB Function so spätnou väzbou stavov}
        \label{fig:mrac_simulink_scheme}
    }
\end{center}

Implementačný kód použitý v bloku MATLAB Function má nasledujúci tvar:

\begin{lstlisting}[language=Matlab,
caption={Implementácia diskrétneho MRAC regulátora v MATLAB Function bloku},
label={lst:mrac_matlab_function}]

function [u, ym, Theta_next, ref_x_next, yf_x_next, rf_x_next, y_prev_next] ...
    = mrac_block(r, y, Theta, ref_x, yf_x, rf_x, y_prev)

% constant parameters
a0m = 1;
a1m = 5;
b0m = 1;

alpha1 = 0.6;
alpha2 = 0.001;
alpha3 = 0.005;

sign_b0 = 1;
Ts = 0.1;

A = [0, 1;
    -a0m, -a1m];

B_ref = [0;
         b0m];

B_f = [0;
       1];

% derivative approximation
ydot = (y - y_prev) / Ts;

% outputs from current states
ym   = ref_x(1);
yf   = yf_x(1);
yf_s = yf_x(2);
rf   = rf_x(1);

% tracking error
e = y - ym;

% control law
omega = [y; ydot; r];
u = Theta.' * omega;

% MIT adaptation law
dTheta = -sign_b0 * e * ...
         [alpha1 * yf;
          alpha2 * yf_s;
          alpha3 * rf];

Theta_next = Theta + Ts * dTheta;

% state updates (Euler method)
ref_x_next = ref_x + Ts * (A * ref_x + B_ref * r);
yf_x_next  = yf_x  + Ts * (A * yf_x  + B_f   * y);
rf_x_next  = rf_x  + Ts * (A * rf_x  + B_f   * r);

% memory update
y_prev_next = y;

end
\end{lstlisting}